\section{实验}

本节利用台湾以东至东海黑潮区域六年的高分辨率海洋再分析数据，按逻辑顺序回答以下五个核心问题：（1）3天前的1°粗分辨率海流图究竟错在哪里，其误差是否具备可由大数据标定的空间物理结构？（2）单纯增加机载计算量能否消除绕障选边的错误，在包住真实海洋的合理集合内环境歧义有多普遍？（3）在所有不取新信息、就地测流、上浮取图与反例门控响应中，哪种策略能真正降低航时损失？（4）航时损失集中在哪些任务上，反例的最坏后悔为何比集合翻转比例更能抓住大损失任务，以及何时才值得付出代价去观测？（5）当海流图变好、航行器变慢或更换误差模型时，本文方法的适用边界在哪里？

\subsection{实验设置与对比基线}
实验严格遵循表\ref{tab:settings}与第4节所述的数据划分：所有误差模态、合理集合标定半径（$r=1.5$）、数值误差容限（$\tau_{\text{num}}=3.2$小时）及门控阈值（$\theta=8$小时）均完全在2019--2021年开发年份（710个任务）上标定选出，随后原封不动地应用于2022--2024年测试年份（623个任务）进行评价。为全面检验计算、决策准则与各类观测手段的效果，我们将全部14种对比方法按其物理机制分为五组（与主对比图一致）：第一组为“同一张粗图上加算”，包括地图+粗算、地图+标准与照地图走（精算）；第二组为“不取新信息、改变决策准则”，包括按期望时间（在32个合理海流样本中选平均航时更短的一侧，即贝叶斯期望准则）与保守（最坏后悔最小，即对所有任务一律选 $\arg\min_s R_s(r)$）；第三组为“就地取新信息”，包括原地停留测（停留2小时测12次）、沿途测30千米、沿途测60千米（边走边测后根据同化更新流场决定是否换边）以及绕行定点测（先绕去前方40千米内反例与地图差异最大处测量）；第四组为“上浮取新图”，包括总是上浮取当天1°图与总是上浮取当天0.5°图（主设置下每次上浮计1小时）；第五组为本文提出的“反例检验分流”，包括门控后取保守、门控后上浮取图（取当天0.5°图）以及作为消融对照的门控后按期望时间。所有结果均以测试年份的平均航时损失（小时）度量，并报告与照地图走（精算）基线配对差值的按天整组自助重抽样（2000次）95\%置信区间。

\subsection{旧粗图错在哪里：黑潮误差的空间结构与尺度}
3天前的1°粗分辨率海流图之所以频繁误导航路决策，根本原因在于其丢失了黑潮主轴及伴生涡旋的高幅值、长相关尺度切变，而这一误差结构可由多年历史再分析数据稳定标定。为了弄清航行器所带粗图的失真规律，我们在开发年份上对比了3天前1°粗图、当天0.5°图与当天1/12°真实海流之间的逐点向量差。图\ref{fig:data}展示了研究海区的平均流态、任务空间分布以及标定所得的误差幅度与空间相关尺度：在图\ref{fig:data}(a)中，黑潮强流带自台湾东北向东北方向蜿蜒贯穿东海陆架坡折，测试任务禁航区广泛覆盖了流轴核心与两侧强切变区；图\ref{fig:data}(b)表明，3天前1°粗图的典型误差并非全区均匀的白噪声，而是沿着黑潮主轴形成一条显著的高误差带，全区典型误差中位数达0.17米/秒（局部流轴处成倍升高）；图\ref{fig:data}(c)进一步显示，3天前1°粗图的误差空间相关长度长达76千米，而当天0.5°图的典型误差中位数则降至0.07米/秒、相关长度缩短至18千米。高达76千米的误差相关长度意味着，当禁航区直径为30千米、起终点跨度为100--140千米时，一个未被粗图解析的中尺度误差模态足以系统性地笼罩整个左侧或右侧通道，使整条备选航线同时遭受逆流或顺流偏置，这为用低维流函数模态构造合理集合提供了坚实的物理依据。

\begin{figure}[htbp]
\centering
\includegraphics[width=\linewidth]{fig_route_data}
\caption{台湾以东至东海研究海区的流场特征与粗图误差标定。(a) 研究海区（23°--31°N、121°--129°E）在2019--2021年开发年份的10米平均海流场，黑潮主轴清晰可见，橙色圆点标示测试任务禁航区中心位置；(b) 3天前1°粗图的典型误差幅度分布（米/秒），最大误差带紧密贴合黑潮强流轴；(c) 海流误差相关系数随两点空间距离的衰减曲线，当天0.5°图的误差相关长度约为18千米，而3天前1°粗图达76千米。旧粗图的误差具有鲜明的空间结构与中尺度相关长度，可由多年海洋大数据可靠标定，构成了构造“合理海流”的物理依据。}
\label{fig:data}
\end{figure}

\subsection{计算消除不了环境歧义：算力与合理反例的实证分析}
实验结果证实，在同一张3天前的1°粗图上成十倍地增加机载计算量只能消除少部分数值歧义，而在能够包住真实海洋的合理集合内，几乎每一个任务都存在翻转决策的合理反例。为了定量分离数值歧义与环境歧义，我们在测试年份的623个任务上分别统计了不同求解精度与不同地图信息下的选错一侧比例，并扫描了合理集合半径 $r$ 从0.25到2.0变化时覆盖真实时间差与存在合理反例的任务占比。如图\ref{fig:ambiguity}(a)所示，当计算量从粗算增加8倍到标准、再增加到64倍的精算时，测试年份选错绕行一侧的任务比例仅从21.7\%微降至20.5\%和19.7\%（在开发年份中从25.1\%降至21.3\%），即增加64倍计算量仅仅修好了粗算错误的28\%，对应的平均航时损失也仅从1.16小时微降至1.03和0.98小时（粗算与精算配对差为+0.18小时，95\%置信区间为$[-0.14, 0.58]$，差异在统计上不显著）；相反，若不增加算力而是将地图换成当天的0.5°图，选错一侧的比例立即减半至10.4\%。图\ref{fig:ambiguity}(b)揭示了多算无效的根本原因：随着合理集合半径 $r$ 取0.25、0.5、1.0、1.5与2.0，测试年份中能找到合理反例的任务比例高达41.6\%、64.8\%、94.1\%、99.7\%与100.0\%；在满足开发年份覆盖率不低于90\%的标定半径 $r=1.5$ 处（开发年份实际包住93\%任务的真实两侧时间差，测试年份包住94\%），高达99.7\%的测试任务都存在至少一个与粗图同样合理、却让另一侧更快的反例海流场。这一结果说明，对于3天前的1°粗图而言，环境歧义是普遍规律而非例外；因此，仅问“是否存在合理反例”并没有区分度，真正决定行动的必然是反例所造成的最坏后悔有多大。

\begin{figure}[htbp]
\centering
\includegraphics[width=\linewidth]{fig_route_ambiguity}
\caption{机载计算量与环境信息对航路决策歧义的不同作用。(a) 测试年份选错绕行一侧的任务比例：在同一张3天前1°粗图上将计算量从粗算（21.7\%）增加至标准（20.5\%）与64倍精算（19.7\%），错误率仅下降约2个百分点，而换成当天0.5°图则直接减半至10.4\%（知道真相为0\%）；(b) 开发年份合理集合包住真实两侧时间差的任务比例与测试年份存在合理反例的任务比例随半径 $r$ 的变化，在标定半径 $r=1.5$ 处两者分别达93\%与99.7\%。机载计算只能消除少量的数值歧义；在足以包住真实海洋的合理集合里，几乎每个决定都存在合理反例，因此必须进一步度量反例导致的最坏损失大小。}
\label{fig:ambiguity}
\end{figure}

\subsection{主对比：五类响应策略的平均航时损失}
在全部五类响应策略的系统对比中，唯有利用合理集合控制尾部风险的保守准则——特别是本文提出的门控后取保守策略——在统计上显著降低了测试年份的平均航时损失，而继续加算、按期望时间选边与各类就地测流均未能改善结局。图\ref{fig:main}汇总了各方法在623个测试任务上的平均损失及置信区间。首先，在不取新信息的方法中，照地图走（精算）的平均损失为0.98小时（95\%置信区间$[0.69, 1.33]$）；改用32个合理海流样本的按期望时间（贝叶斯）准则后，平均损失反而略升至1.04小时（与精算地图配对差为+0.06小时，$[-0.22, 0.34]$），即使加上反例门控（门控后按期望时间）平均损失仍为0.96小时（配对差-0.01小时，$[-0.29, 0.25]$），说明在对称分布的随机扰动下求平均航时并不能挑出翻转决策的不利流场。相反，本文提出的门控后取保守策略将测试年份的平均损失大幅降至0.70小时（$[0.48, 0.97]$），与照地图走（精算）的配对差为-0.28小时，95\%置信区间为$[-0.51, -0.07]$（严格不包含0），相对损失降低了28\%（在开发年份中同样从精算地图的1.13小时降至0.70小时）。同时可以如实看到，对所有任务一律取保守（最坏后悔最小）的平均损失为0.72小时（与精算地图配对差-0.26小时，$[-0.49, -0.05]$；开发年份为0.76小时），门控后取保守比一律取保守仅再降0.02小时：这表明在无需付出额外观测代价时，核心收益来自反例搜索所算出的最坏后悔准则本身，而门控的作用在于精准锁定了那47\%的高风险任务（其间保守规则实际改选另一侧的任务仅41个），并在需要调用有代价的上浮观测时避免在其余53\%的安全任务上浪费时间（总是上浮取当天0.5°图损失为1.22小时，而门控后上浮取0.5°图降至0.78小时，开发年份为0.76小时）。

\begin{figure}[htbp]
\centering
\includegraphics[width=\linewidth]{fig_route_main}
\caption{测试年份623个强流任务上五类响应策略的平均航时损失对比（竖直虚线为照精算地图走的平均损失0.98小时）。各组从上至下依次为：同一张粗图上加算（粗算1.16、标准1.03、精算0.98小时）；不取新信息而改变决策准则（按期望时间1.04、最坏后悔最小的保守准则0.72小时）；就地取新信息（原地停留测2.91、沿途测30千米0.93、沿途测60千米0.91、绕行定点测12.43小时）；总是上浮取新图（每次计1小时：当天1°图1.67、当天0.5°图1.22小时）；以及本文反例检验门控分流（门控后取保守0.70、门控后上浮取0.5°图0.78小时）。只有基于合理反例最坏后悔的保守与门控策略显著降低了损失，单纯加算与就地测流均难以奏效。}
\label{fig:main}
\end{figure}

\subsection{损失集中在哪里、何时值得观测}
航路决策的航时损失在测试任务中呈现极端的长尾集中分布，反例的最坏后悔之所以有效，正是因为它比集合翻转比例和地图账面时间差更能将那些一旦选错就会损失十几个小时的高危任务排在最前面。统计表明，照地图走（精算）损失最大的前10\%任务竟然占了测试年份全部航时损失的92\%（开发年份占88\%）。如果只看一个指标能否区分“是否选错且多花1小时以上”（测试年份共79个此类任务）而不论具体多花多少小时，反例最坏后悔、32成员集合翻转比例与地图两侧时间差三者的受试者工作特征曲线下面积其实十分接近（分别为0.65、0.67与0.67）。然而，一旦考察谁能抓住真正造成巨大航时损失的任务，三者便拉开显著差距：如图\ref{fig:risk}(a)所示，按风险指标从高到低排序时，排在前20\%的任务所捕获的地图总损失比例，按反例最坏后悔高达67\%，而按32个合理场的翻转比例仅为49\%，按地图上两侧时间差越小越危险排则只有34\%（随机排序为20\%，按事后真实损失排的理论上限为100\%）；在门控阈值 $\theta=8$ 小时下，被标记的47\%任务更是囊括了高达82\%的地图总损失。这是因为集合翻转比例对“翻转后多花0.5小时”与“翻转后多花15小时”一视同仁，而梯度搜索的反例最坏后悔直接度量了不利流场在两侧通道造成的流速积分差。

进一步分析各类观测手段为何有效或失效，可以清楚看到观测的空间覆盖范围与时间代价之间的硬约束。首先，就地原位测流在全部测试任务上均未能带来显著净收益：出发前原地停留测平均损失高达2.91小时（配对差+1.94小时，$[1.55, 2.31]$），因为它白白耗去了2小时驻留时间，却只在起点一个点上消除了0.05米/秒的仪器测量噪声，对远方禁航区两侧的海流毫无感知；专门绕行到前方40千米内反例与地图差异最大处去定点测量更是灾难性的，平均损失高达12.43小时（配对差+11.45小时，$[10.49, 12.46]$），因为航行器平均要先花15.4小时才能顶流赶到该测量点；而顺路沿所选一侧航行30千米或60千米边走边测后允许换边，平均损失分别为0.93小时（配对差-0.05小时，$[-0.30, 0.18]$）与0.91小时（配对差-0.07小时，$[-0.27, 0.05]$），改善均不显著——具体而言，沿途测30千米后决定调头换边的任务仅占3.4\%（其中81\%换对），因为此时尚未进入禁航区两侧的核心切变带，同化更新后的流场几乎不改变原判断；而走到60千米再换边的任务虽增至18.8\%，但由于已在原侧深入近半程，此时调头绕回另一侧的折返代价过大，换边反而只有15\%是划算的。其次，对于能够一眼看清全局的“上浮取当天0.5°图”，图\ref{fig:risk}(b)展示了平均损失随单次上浮代价（0至4小时）的变化规律：当上浮代价为0、0.5、1.0、2.0与3.0小时时，总是上浮取0.5°图的损失随代价线性攀升（分别为0.22、0.72、1.22、2.22与3.22小时），而门控后上浮取0.5°图因为只在47\%的高危任务上触发上浮，损失分别控制在0.22、0.55、0.78、1.06与0.98小时。由此得出清晰的行动边界：当单次上浮代价低于约1小时时，门控后上浮取新图是最优响应；当上浮代价高于约1小时时，无需观测的门控后取保守（恒为0.70小时）成为最优选择。

\begin{figure}[htbp]
\centering
\includegraphics[width=\linewidth]{fig_route_risk}
\caption{航时损失的风险集中度排序与上浮观测代价平衡分析。(a) 按不同风险指标从高到低排序时，前一定比例任务累计捕获的地图总损失占比：排名前20\%的任务中，反例最坏后悔捕获了67\%的损失，显著高于32个合理场翻转比例（49\%）与地图两侧时间差（34\%）；(b) 测试年份平均损失随单次上浮取当天0.5°图时间代价（0--4小时）的变化曲线：总是上浮随代价线性恶化，门控后上浮在代价低于约1小时时最优，代价高于约1小时后门控后取保守（0.70小时）最优。反例最坏后悔能够精准揪出大损失任务，从而实现“观测便宜时门控去测、观测昂贵时门控保守”的自适应切换。}
\label{fig:risk}
\end{figure}

\subsection{适用范围与消融分析}
为了检验合理反例检验在不同先验地图质量、不同航行器机动能力以及不同合理集合构造方法下的适用边界，我们在另外三组设置（每组测试年份约350个任务，门控阈值各自在开发年份按相同规则标定）上开展了对比与消融实验，结果汇总于表\ref{tab:scope}。第一，当航行器下水时直接带有质量较好的当天0.5°图时（测试349个任务，平均最快航时38.7小时，标定半径 $r=1.75$），由于背景误差大幅减小，存在合理反例的任务比例从99.7\%降至64.5\%，精算地图选错率仅8.0\%（粗算8.9\%），照精算地图走的平均损失已低至0.17小时；此时反例最坏后悔前20\%任务仍能集中60\%的损失，且门控后取保守（0.17小时，与地图配对差+0.00小时，$[-0.04, 0.05]$）与门控后上浮取图（0.16小时）自动退化为绝大多数时候照地图走，不会产生负向干扰。第二，当航行器为对水速度 $V=0.75$ 米/秒的慢速航行器时（测试347个任务，平均最快航时延长至59.8小时，标定半径 $r=1.5$），流速接近甚至超过航速的风险加剧，虽然精算地图选错率为18.4\%（粗算22.2\%，存在合理反例占99.1\%），但一旦选错付出的航时成倍放大，照精算地图走的平均损失飙升至2.43小时；此时反例最坏后悔前20\%任务集中了高达74\%的损失，门控后取保守将平均损失降至1.80小时（减少0.63小时，但95\%区间$[-1.67, 0.30]$含0，未达显著），而由于选错的代价变大、1小时的上浮代价相对变小，门控后上浮取0.5°图将损失降至1.32小时。第三，在消融实验中，若去掉本文用多年数据标定的流函数无散误差模态，改用从历史相似日直接抽取小尺度差值场来拼装合理集合（测试349个任务，平均最快航时39.8小时，标定半径 $r=1.5$，精算选错19.5\%、粗算22.3\%，存在反例99.4\%），虽然前20\%任务仍能捕获64\%损失且门控后上浮取图为0.77小时，但门控后取保守的平均损失为0.95小时，相比精算地图的0.99小时仅微降0.04小时（$[-0.43, 0.36]$），几乎失去了保守选边的纠错能力。两组的反例搜索方法完全相同，差别只在“合理”如何定义：由多年数据标定误差大小与相关长度的合理集合，能够指出真正危险的一侧，而拼装历史相似日的小尺度差则不能。

\begin{table}[htbp]
\centering
\caption{不同海流图质量、航行器速度及合理集合构造方法下的适用范围与消融实验结果（损失单位：小时）。}
\label{tab:scope}
\small
\begin{tabular}{@{}p{0.19\linewidth}>{\centering\arraybackslash}p{0.088\linewidth}>{\centering\arraybackslash}p{0.088\linewidth}>{\centering\arraybackslash}p{0.088\linewidth}>{\centering\arraybackslash}p{0.088\linewidth}>{\centering\arraybackslash}p{0.088\linewidth}>{\centering\arraybackslash}p{0.088\linewidth}>{\centering\arraybackslash}p{0.088\linewidth}@{}}
\toprule
设置 & 测试任务 & 存在合理反例 & 基线选错率 & 基线额外航时 & 超阈值时保守改选 & 超阈值时上浮取图 & 最大后悔值前20\%所占额外航时 \\
\midrule
黑潮主设置 & 623 & 99.7\% & 19.7\% & 0.98 & 0.70 & 0.78 & 67\% \\
粗图换成当天0.5°图 & 349 & 64.5\% & 8.0\% & 0.17 & 0.17 & 0.16 & 60\% \\
航速降至0.75米/秒 & 347 & 99.1\% & 18.4\% & 2.43 & 1.80 & 1.32 & 74\% \\
消融：合理集合改用相似日差值 & 349 & 99.4\% & 19.5\% & 0.99 & 0.95 & 0.77 & 64\% \\
湾流 & 643 & 98.8\% & 19.4\% & 3.56 & 4.00 & 1.31 & 71\% \\
阿古拉斯回流区 & 618 & 99.2\% & 27.7\% & 6.87 & 6.49 & 2.20 & 53\% \\
改带卫星测高地转流图 & 648 & 81.0\% & 19.6\% & 0.79 & 0.71 & 0.64 & 67\% \\
\bottomrule
\end{tabular}

\end{table}

\subsection{局限性讨论}
本文的研究与实验设计存在四个方面的局限。第一，实验在求取单次任务航线时将每张海流快照视为空间变化但时间定常的二维流场，尚未在约40小时的航行过程中耦合流场随时间的连续演变。第二，实验以当天1/12°的再分析海流场作为检验决策的“真实海洋”，尽管其很好地解析了黑潮主轴与中尺度涡旋，但仍平滑了数千米以下的亚中尺度过程与内波扰动。第三，本文将航路规划简化为绕越单个圆形禁航区的左、右二元拓扑决策，而在包含多岛礁群的复杂海域中，备选拓扑通道将扩展为多分支同伦类。第四，实验假设航行器在选定左或右一侧后，能够依靠局部走航感知在该侧通道内按真实海流以最优轨迹航行，这一理想化执行假设使本文得以剥离底层轨迹跟踪控制误差，纯粹评估顶层二元选边决策的得失。
